% !TEX root =  Lecture10.tex


%%%%%%%%%%%%%%%%%%%%%%%%%%%%%%%%%%%%
% Recap/Agenda
%%%%%%%%%%%%%%%%%%%%%%%%%%%%%%%%%%%%
\begin{frame}
\frametitle{Module 2: Statistical Inference}
    \begin{itemize}
    	\item \hl{Previously:} testing with randomization --- \emph{is there an effect?}
        \item \hl{This time:} \emph{how big is it?} Confidence intervals with bootstrapping.
          \begin{itemize}
              \item confidence intervals as statements about a \emph{procedure};
              \item the bootstrap idea: resample your own data with replacement;
              \item percentile and standard-error intervals, and what makes them narrow;
              \item doing it in R --- and what the bootstrap cannot do.
          \end{itemize}
        \item \hl{Reading:} IMS Chapter 12.
        \item \hl{Homework for today:}
        \item \hl{Homework for next class:}
    \end{itemize}
\end{frame}


%%%%%%%%%%%%%%%%%%%%%%%%%%%%%%%%%%%%
% Learning objectives:
%%%%%%%%%%%%%%%%%%%%%%%%%%%%%%%%%%%%
\begin{frame}
    \frametitle{Lecture Objectives}
    \goals{%
    \begin{itemize}
    \item say what population \emph{parameter} a scientific question is really asking about;
    \item build a bootstrap distribution --- resample, recompute, repeat --- and read a confidence interval off it, by percentile or by standard error;
    \item interpret a confidence interval correctly, and name the three standard misreadings;
    \item carry the whole workflow out in R, and choose between randomization, bootstrapping, and a mathematical model.
    \end{itemize}}
    \objectives{%
    \begin{itemize}
    \item \textbf{(M2, LO1)} Approaches to Inference
    \item \textbf{(M2, LO2)} Simulation-based Inference
    \end{itemize}}
\end{frame}
